% concatenated from document.tex
% !TEX program = pdflatex
\documentclass[11pt,a4paper]{article}

% ---------- Basic encoding and fonts ----------
\usepackage[utf8]{inputenc}       % Input encoding
\usepackage[T1]{fontenc}          % Output encoding
\usepackage{microtype}            % Fine-tune typesetting

% ---------- Mathematical symbols and theorem environments ----------
\usepackage{amsmath, amssymb}     % American Mathematical Society packages
\usepackage{amsthm}               % Theorem environments
\usepackage{mathtools}            % Extended amsmath
\usepackage{bm}                   % Bold math symbols
\usepackage{comment}

% Define theorem-like environments (add or remove as needed)
\theoremstyle{plain}              % Body text in italics
\newtheorem{theorem}{Theorem}[section]
\newtheorem{lemma}[theorem]{Lemma}
\newtheorem{corollary}[theorem]{Corollary}
\newtheorem{proposition}[theorem]{Proposition}

\theoremstyle{definition}         % Body text in upright
\newtheorem{definition}[theorem]{Definition}
\newtheorem{example}[theorem]{Example}
\newtheorem{remark}[theorem]{Remark}
\newtheorem{note}[theorem]{Note}

% ---------- Figures and floats ----------
\usepackage{graphicx}             % Insert figures
\usepackage{float}                % Provides the H option
\usepackage[font=small,labelfont=bf]{caption} % Captions for figures and tables
\usepackage{subcaption}           % Subfigures

% ---------- Pseudocode ----------
\usepackage{algorithm}            % Algorithm float
\usepackage{algpseudocode}        % Algorithm style (requires algorithmicx)
% Replace with algorithm2e if preferred; the more general algpseudocode is used here

% ---------- Page setup ----------
\usepackage[margin=2.5cm]{geometry}    % Page margins
\usepackage{parskip}% Leave a blank line between paragraphs; indentation is handled by parskip
\setlength{\parindent}{0pt}% Cancel first-line indentation (optional, based on personal preference)

% ---------- Cross-references and hyperlinks ----------
\usepackage{xcolor} % Optional
\usepackage[
colorlinks=true,
linkcolor=blue,% Internal cross-references: \ref, \cref, \autoref, etc.
citecolor=green,% Bibliography citations: \cite, etc.
urlcolor=red% URLs: \url, \href, etc.
]{hyperref}
\usepackage[nameinlink]{cleveref} % Still must come after hyperref

% Set reference format (optional)
\crefname{equation}{equation}{equations}
\crefname{figure}{figure}{figures}
\crefname{table}{table}{tables}
\crefname{algorithm}{algorithm}{algorithms}

\title{\textbf{Element-wise Convergence Behavior of Subspace Iteration}}
\author{Mingyang Zhao\\
myzhao24@m.fudan.edu.cn
}
\date{\today}

\begin{document}

\maketitle

%\tableofcontents

\begin{abstract}
This paper studies the element-wise convergence behavior of subspace iteration. All results are established for both \(\mathbb{F}=\mathbb{R}\) and \(\mathbb{F}=\mathbb{C}\). For a diagonal matrix \(\Lambda=\mathrm{diag}(\lambda_1,\ldots,\lambda_n)\) with \(|\lambda_1|>\cdots>|\lambda_n|>0\), we analyze the modulus of each entry of the iterative matrix sequence generated by subspace iteration, both without and with the Rayleigh--Ritz procedure. Under mild assumptions on the initial matrix \(X\in\mathbb{F}^{n\times m}\), we first derive exact asymptotic expressions for \(|Q_k(i,j)|\) in the subspace iteration without the Rayleigh--Ritz procedure: entries with \(i\neq j\) decay as \((|\lambda_{\max\{i,j\}}|/|\lambda_{\min\{i,j\}}|)^k\), and the deviation of \(|Q_k(j,j)|\) from \(1\) decays as \(\max\{|\lambda_j|/|\lambda_{j-1}|,|\lambda_{j+1}|/|\lambda_j|\}^{2k}\), with explicit coefficients determined by the LU factorization of \(X\). For the Rayleigh--Ritz variant, we obtain element-wise bounds for the normalized Ritz vectors \(Z_k\). Specifically, for \(|Z_k(i,j)|\), off-diagonal entries with \(i\geq m+1\) decay as \((|\lambda_i|/|\lambda_j|)^k\), off-diagonal entries with \(i\leq m\) decay as \((|\lambda_{m+1}|^2/(|\lambda_i||\lambda_j|))^k\), and the deviation of \(|Z_k(j,j)|\) from \(1\) decays as \((|\lambda_{m+1}|/|\lambda_j|)^{2k}\). These results give an element-wise description of the convergence behavior of subspace iteration. After an orthogonal or unitary change of basis, the results apply to real symmetric or complex normal matrices.
\end{abstract}

\textbf{Keywords:} subspace iteration, Rayleigh--Ritz procedure.

\section{Introduction}

Subspace iteration is a fundamental iterative method for computing a dominant invariant subspace of a large square matrix. It can be viewed as a block generalization of the power method: instead of iterating a single vector, one propagates an \(n\times m\) block and orthonormalizes its columns at each step. Because of its simple structure and its close relation to the QR algorithm and Krylov subspace methods, it remains a standard tool for large-scale eigenvalue problems and has been widely used in scientific computing and data analysis; see, e.g., \cite{golub_book,parlett_book,saad_book}.

The pseudocode of subspace iteration without the Rayleigh--Ritz procedure is as follows \cite{saad_book}.

\begin{algorithm}[H]
\caption{Subspace iteration without the Rayleigh--Ritz procedure}
\begin{algorithmic}[1]
\Require \(A\in\mathbb{C}^{n\times n}\), \(X\in\mathbb{C}^{n\times m}\)
\State \(Q_0R_0=X\)
\For{\(k=1,2,\dots\)}
\State \(Q_kR_k=AQ_{k-1}\)
\EndFor
\end{algorithmic}
\end{algorithm}

For simplicity, assume that \(A=\mathrm{diag}(\lambda_1,\dots,\lambda_n)\) with \(|\lambda_1|>\dots>|\lambda_n|>0\). Under mild assumptions on \(X\), we have \(|Q_k|\to\begin{bmatrix}
	I_m\\
	O
\end{bmatrix}\), that is,
\[
|Q_k(i,j)|\to\delta_{i,j},\quad (i,j)\in[n]\times[m].
\]
However, what is the convergence behavior of the modulus of each entry \(|Q_k(i,j)|\)?

The pseudocode of subspace iteration with the Rayleigh--Ritz procedure is as follows \cite{saad_book}.

\begin{algorithm}[H]
\caption{Subspace iteration with the Rayleigh--Ritz procedure}
\begin{algorithmic}[1]
\Require \(A\in\mathbb{C}^{n\times n}\), \(X\in\mathbb{C}^{n\times m}\)
\State \(Q_0R_0=X\)
\State \(Q_0^*AQ_0=U_0\Theta_0U_0^{-1}\)
\State \(Z_0=Q_0U_0\)
\For{\(k=1,2,\dots\)}
\State \(Q_kR_k=AZ_{k-1}\)
\State \(Q_k^*AQ_k=U_k\Theta_kU_k^{-1}\)
\State \(Z_k=Q_kU_k\)
\EndFor
\end{algorithmic}
\end{algorithm}

For simplicity, assume that \(A=\mathrm{diag}(\lambda_1,\dots,\lambda_n)\) with \(|\lambda_1|>\dots>|\lambda_n|>0\) Under mild assumptions on \(X\), we have \(|Z_k|\to\begin{bmatrix}
	I_m\\
	O
\end{bmatrix}\), that is,
\[
|Z_k(i,j)|\to\delta_{i,j},\quad (i,j)\in[n]\times[m].
\]
However, what is the convergence behavior of the modulus of each entry \(|Z_k(i,j)|\)?

Classical convergence theory for subspace iteration is usually formulated either at the subspace level or at the level of individual eigenvectors and Ritz vectors. It is worth emphasizing that these classical results are not element-wise: they do not describe the decay or the deviation of individual entries of the computed basis matrices \(Q_k\) or \(Z_k\).

At the subspace level, Theorem~7.3.1 in \cite{golub_book} analyzes orthogonal iteration, which is equivalent to subspace iteration in the absence of the Rayleigh--Ritz procedure. It shows that the distance between the computed subspace \(\operatorname{ran}(Q_k)\) and the dominant invariant subspace \(D_r(A)\) is bounded by a constant times \((|\lambda_{r+1}|/|\lambda_r|)^k\). The constant depends on the separation of the dominant eigenvalues and on the departure from normality of \(A\). This is a subspace-level convergence result: it measures the convergence of the whole computed subspace, not the convergence of its individual entries.

At the vector level, Theorem~14.4.1 in \cite{parlett_book} gives a convergence estimate for each target eigenvector \(z_i\) in the symmetric case. More precisely, it shows that a certain vector \(x_i^{(k)}\) constructed from the iterated subspace approximates \(z_i\) linearly, with a rate governed by ratios of the relevant eigenvalues. Theorem~14.4.2 in \cite{parlett_book} further shows that the Ritz vector \(y_i^{(k)}\) is asymptotically equivalent to \(x_i^{(k)}\), and therefore the Ritz vector itself converges to the corresponding eigenvector at the same vector-level rate. These are vector-level results for individual eigenvectors and Ritz vectors.

Also at the vector level, Theorem~5.2 in \cite{saad_book} considers subspace iteration with projection. For each wanted eigenvector \(u_i\), it estimates the distance from \(u_i\) to the iterated subspace \(S_k\). The bound decays like \((|\lambda_{m+1}|/|\lambda_i|)^k\), up to a factor that tends to zero as \(k\to\infty\). This is again a per-eigenvector estimate: it describes how fast each wanted eigenvector is captured by the iterated subspace, rather than how the subspace as a whole converges or how individual entries of the basis behave.

In summary, the existing theory provides either subspace-level bounds, as in Theorem~7.3.1 of \cite{golub_book}, or vector-level bounds, as in Theorems~14.4.1 and~14.4.2 of \cite{parlett_book} and Theorem~5.2 of \cite{saad_book}. None of these results gives an element-wise description of the convergence of subspace iteration. The purpose of this paper is to fill this gap by analyzing the modulus of each entry of the iterative matrices, both without and with the Rayleigh--Ritz procedure.

The main contributions of this paper are as follows.
\begin{itemize}
\item For subspace iteration without the Rayleigh--Ritz procedure, we derive exact asymptotic expressions for the modulus of every entry of the orthonormal basis \(Q_k\). For \(i\ne j\), the entry \(|Q_k(i,j)|\) decays as
\[
\left(\frac{|\lambda_{\max\{i,j\}}|}{|\lambda_{\min\{i,j\}}|}\right)^k,
\]
while the deviation of \(|Q_k(j,j)|\) from \(1\) decays as
\[
\max\left\{\frac{|\lambda_j|}{|\lambda_{j-1}|},
\frac{|\lambda_{j+1}|}{|\lambda_j|}\right\}^{2k}.
\]
The coefficients in these asymptotic expressions are explicitly determined by the LU factorization of the initial matrix \(X\).
\item For subspace iteration with the Rayleigh--Ritz procedure, we establish element-wise upper and lower bounds for the normalized Ritz vectors \(Z_k\). In particular, off-diagonal entries with \(i\ge m+1\) decay as
\[
\left(\frac{|\lambda_i|}{|\lambda_j|}\right)^k,
\]
off-diagonal entries with \(i\le m\) decay as
\[
\left(\frac{|\lambda_{m+1}|^2}{|\lambda_i|\,|\lambda_j|}\right)^k,
\]
and the deviation of \(|Z_k(j,j)|\) from \(1\) decays as
\[
\left(\frac{|\lambda_{m+1}|}{|\lambda_j|}\right)^{2k}.
\]
\item All results are proved for both \(\mathbb{F}=\mathbb{R}\) and \(\mathbb{F}=\mathbb{C}\). After an orthogonal or unitary change of basis, they apply to real symmetric or complex normal matrices.
\end{itemize}

The proofs use two different technical routes. In the absence of the Rayleigh--Ritz procedure, we exploit the LU factorization of the initial block and perform a Gram--Schmidt orthogonalization of the columns of the scaled lower triangular factor; this yields an inductive element-wise description of \(Q_k\). For the Rayleigh--Ritz variant, the analysis is more delicate. We reformulate the Ritz pair equations as a fixed point problem in a suitably weighted metric space and then apply the Banach fixed point theorem to obtain quantitative bounds for the Ritz values and the coefficients of the Ritz vectors.

Although the entrywise convergence of subspace iteration has also been studied in \cite{NEURIPS2020_3d8e03e8}, our work is fundamentally different. The analysis in \cite{NEURIPS2020_3d8e03e8} measures the error in the \(\ell_{2\to\infty}\) norm and provides only upper bounds, whereas we establish element-wise asymptotic results: exact expressions with explicit coefficients for the case without the Rayleigh-Ritz procedure, and element-wise upper and lower bounds for the Rayleigh-Ritz variant. This difference in both the convergence metric and the precision of the results makes the two analyses essentially distinct.

The rest of this paper is organized as follows. Section~\ref{section:withoutRR} is devoted to subspace iteration without the Rayleigh--Ritz procedure. We first recall the LU factorization of the initial block and the associated scaled lower triangular matrices, and then prove the exact asymptotic element-wise formulas by induction. Section~\ref{section:withRR} treats the Rayleigh--Ritz variant. After introducing the necessary notation and auxiliary lemmas, we construct a complete metric space and a contraction map, apply the Banach fixed point theorem, and finally translate the fixed point estimates back to the normalized Ritz vectors. Section~\ref{section:conclusion} concludes the paper and discusses a possible extension to unitarily upper triangularizable matrices.

\section{Subspace iteration without the Rayleigh--Ritz procedure}\label{section:withoutRR}

\begin{lemma}\label{lemma:withoutRR}
Suppose \(\mathbb{F}=\mathbb{R}\) or \(\mathbb{C}\), \(n\), \(m\in\mathbb{N}_+\) with \(m\leq n\), and \(\Lambda=\mathrm{diag}(\lambda_1,\dots,\lambda_n)\in\mathbb{F}^{n\times n}\), with \(|\lambda_1|>\dots>|\lambda_n|>0\). Assume \(X\in\mathbb{F}^{n\times m}\) satisfies \(\det(X(1:j,1:j))\neq 0\) for every \(j\in[m]\). Then there exists a unique unit lower triangular \(L\in\mathbb{F}^{n\times m}\) and upper triangular \(U\in\mathbb{F}^{m\times m}\) such that \(X=LU\). Assume further that all strictly lower triangular entries of \(L\) and of \(L(1:m,:)^{-1}\) are nonzero.

Let \(Q_kR_k=\Lambda^k X\) be a QR factorization of \(\Lambda^kX\). Then for every \(j\in[m]\), we have
\[
|Q_k(i,j)|=\begin{cases}
|(L(1:m,:)^{-1})(j,i)|\dfrac{|\lambda_j|^k}{|\lambda_i|^k}(1+o(1)),\quad&i\in[j-1]\\
1+o(1),\quad&i=j\\
|L(i,j)|\dfrac{|\lambda_i|^k}{|\lambda_j|^k}(1+o(1)),\quad&i\in[n]\setminus[j]
\end{cases}.
\]
\end{lemma}

\begin{theorem}
Suppose \(\mathbb{F}=\mathbb{R}\) or \(\mathbb{C}\), \(n\), \(m\in\mathbb{N}_+\) with \(m\leq n\) and \(\Lambda=\mathrm{diag}(\lambda_1,\dots,\lambda_n)\in\mathbb{F}^{n\times n}\), with \(|\lambda_1|>\dots>|\lambda_n|>0\). Suppose \(V\in\mathbb{F}^{n\times n}\) satisfies \(VV^*=V^*V=I_n\) and define \(A:=V\Lambda V^*\). Suppose \(Y\in\mathbb{F}^{n\times m}\) and define \(X:=V^*Y\). Assume \(\det(X(1:j,1:j))\neq 0\) for every \(j\in[m]\). Then there exists a unique unit lower triangular \(L\in\mathbb{F}^{n\times m}\) and upper triangular \(U\in\mathbb{F}^{m\times m}\) such that \(X=LU\). Suppose further that all strictly lower triangular entries of \(L\) and of \(L(1:m,:)^{-1}\) are nonzero.

Let \(\tilde{Q}_k\tilde{R}_k=A^kY\) be a QR factorization of \(A^kY\). Then for every \(j\in[m]\), we have
\[
|V(:,i)^*\tilde{Q}_k(:,j)|=\begin{cases}
|(L(1:m,:)^{-1})(j,i)|\dfrac{|\lambda_j|^k}{|\lambda_i|^k}(1+o(1)),\quad&i\in[j-1]\\
1+o(1),\quad&i=j\\
|L(i,j)|\dfrac{|\lambda_i|^k}{|\lambda_j|^k}(1+o(1)),\quad&i\in[n]\setminus[j]
\end{cases}.
\]

In particular, we have
\[
|V(:,i)^*\tilde{Q}_k(:,j)|=\Theta\left(\dfrac{|\lambda_{\max\{i,j\}}|^k}{|\lambda_{\min\{i,j\}}|^k}\right),\quad i\neq j;
\]
\[
1-|V(:,j)^*\tilde{Q}_k(:,j)|=\Theta\left(\max\left\{\dfrac{|\lambda_j|}{|\lambda_{j-1}|},\dfrac{|\lambda_{j+1}|}{|\lambda_j|}\right\}^{2k}\right).
\]
\end{theorem}
\begin{proof}
This follows by applying a change of orthonormal basis to the conclusion of the preceding lemma.
\end{proof}

We now prove Lemma~\ref{lemma:withoutRR}.

\subsection{LU factorization of \(X\)}

Since \(\det(X(1:j,1:j))\neq 0\) for all \(j\in[m]\), there exists a unique unit lower triangular \(L\in\mathbb{F}^{n\times m}\) and upper triangular \(U\in\mathbb{F}^{m\times m}\) such that \(X=LU\). Let \(\Lambda_1=\Lambda(1:m,1:m)\). Set \(L_k=\Lambda^kL\Lambda_1^{-k}\). Then \(L_k\in\mathbb{F}^{n\times m}\) remains unit lower triangular, and its strictly lower triangular entries are:
\[
L_k(i,j)=L(i,j)\left(\dfrac{\lambda_i}{\lambda_j}\right)^k,\quad\forall i>j.
\]

\subsection{Gram--Schmidt orthogonalization of the columns of \(L_k\)}

For \(j=1,2,\dots,m\), define \(V_k(:,j)\) and \(Q_k(:,j)\) inductively by
\[
V_k(:,j)=L_k(:,j)-\sum_{l=1}^{j-1}\langle L_k(:,j),Q_k(:,l)\rangle Q_k(:,l),\quad Q_k(:,j)=\dfrac{V_k(:,j)}{\|V_k(:,j)\|}.
\]

Claim: For every \(j\in[m]\), we have
\[
Q_k(i,j)=\begin{cases}
\overline{(L(1:m,:)^{-1})(j,i)}\overline{\left(\dfrac{\lambda_j}{\lambda_i}\right)^k}(1+o(1)),\quad&i\in[j-1]\\
1+o(1),\quad&i=j\\
L(i,j)\left(\dfrac{\lambda_i}{\lambda_j}\right)^k(1+o(1)),\quad&i\in[n]\setminus[j]
\end{cases}.
\]

We prove the claim by induction.

\subsection{Base case}

Since \(\|V_k(:,1)\|_2=\|L_k(:,1)\|_2=1+o(1)\), for every \(i\in[n]\) we have
\[
Q_k(i,1)=V_k(i,1)/\|V_k(:,1)\|=L_k(i,1)(1+o(1))=L(i,1)\left(\dfrac{\lambda_i}{\lambda_1}\right)^k(1+o(1)).
\]

\subsection{Inductive step}

Assume the claim holds for \(1\), \(2\), \(\dots\), \(j-1\). We prove it for \(j\).

For \(l\in[j-1]\), define \(c_{l,j}=\langle L_k(:,j), Q_k(:,l)\rangle\). Then
\begin{align*}
c_{l,j}&=\sum_{i=j}^{n}L_k(i,j)\overline{Q_k(i,l)}\\
&=\sum_{i=j}^{n}L(i,j)\left(\dfrac{\lambda_i}{\lambda_j}\right)^k\overline{L(i,l)\left(\dfrac{\lambda_i}{\lambda_l}\right)^k}(1+o(1))\\
&=\overline{L(j,l)\left(\dfrac{\lambda_j}{\lambda_l}\right)^k}(1+o(1)).
\end{align*}

\subsubsection{Computation of \(V_k(i,j)\) (\(i>j\))}

\begin{align*}
V_k(i,j)&=L_k(i,j)-\sum_{l=1}^{j-1}c_{l,j}Q_k(i,l)\\
&=L(i,j)\left(\dfrac{\lambda_i}{\lambda_j}\right)^k-\sum_{l=1}^{j-1}\overline{L(j,l)\left(\dfrac{\lambda_j}{\lambda_l}\right)^k}(1+o(1))L(i,l)\left(\dfrac{\lambda_i}{\lambda_l}\right)^k(1+o(1))\\
&=L(i,j)\left(\dfrac{\lambda_i}{\lambda_j}\right)^k(1+o(1)).
\end{align*}

\subsubsection{Computation of \(V_k(j,j)\)}

\begin{align*}
V_k(j,j)&=L_k(j,j)-\sum_{l=1}^{j-1}c_{l,j}Q_k(j,l)\\
&=1-\sum_{l=1}^{j-1}\overline{L(j,l)\left(\dfrac{\lambda_j}{\lambda_l}\right)^k}(1+o(1))L(j,l)\left(\dfrac{\lambda_j}{\lambda_l}\right)^k(1+o(1))\\
&=1+o(1).
\end{align*}

\subsubsection{Computation of \(V_k(i,j)\) (\(i<j\))}

\begin{align*}
V_k(i,j)&=L_k(i,j)-\sum_{l=1}^{j-1}c_{l,j}Q_k(i,l)\\
&=-\sum_{l=1}^{i-1}c_{l,j}Q_k(i,l)-\sum_{l=i}c_{l,j}Q_k(i,l)-\sum_{l=i+1}^{j-1}c_{l,j}Q_k(i,l)\\
&=-\sum_{l=1}^{i-1}\overline{L(j,l)\left(\dfrac{\lambda_j}{\lambda_l}\right)^k}(1+o(1))L(i,l)\left(\dfrac{\lambda_i}{\lambda_l}\right)^k(1+o(1))\\
&\quad-\overline{L(j,i)\left(\dfrac{\lambda_j}{\lambda_i}\right)^k}(1+o(1))(1+o(1))\\
&\quad-\sum_{l=i+1}^{j-1}\overline{L(j,l)\left(\dfrac{\lambda_j}{\lambda_l}\right)^k}(1+o(1))\overline{(L(1:m,:)^{-1})(l,i)\left(\dfrac{\lambda_l}{\lambda_i}\right)^k}(1+o(1))\\
&=\left(-\overline{L(j,i)}-\sum_{l=i+1}^{j-1}\overline{L(j,l)(L(1:m,:)^{-1})(l,i)}\right)\overline{\left(\dfrac{\lambda_j}{\lambda_i}\right)^k}(1+o(1)).
\end{align*}

Since \(L(1:m,:)\) and \(L(1:m,:)^{-1}\) are both unit lower triangular matrices, we have
\begin{align*}
0&=(L(1:m,:)L(1:m,:)^{-1})(j,i)=\sum_{l=i}^{j}L(j,l)L(1:m,:)^{-1}(l,i)\\
&=L(j,i)+\sum_{l=i+1}^{j-1}L(j,l)L(1:m,:)^{-1}(l,i)+L(1:m,:)^{-1}(j,i),
\end{align*}
so
\[
\overline{L(1:m,:)^{-1}(j,i)}=-\overline{L(j,i)}-\sum_{l=i+1}^{j-1}\overline{L(j,l)L(1:m,:)^{-1}(l,i)}.
\]
Hence
\[
V_k(i,j)=\overline{L(1:m,:)^{-1}(j,i)}\overline{\left(\dfrac{\lambda_j}{\lambda_i}\right)^k}(1+o(1)).
\]

\subsubsection{Normalization in the 2-norm}

Since \(\|V_k(:,j)\|_2=1+o(1)\), we have \(Q_k(:,j)=V_k(:,j)(1+o(1))\), hence the claim holds for \(j\).

\subsection{Completing the proof}

By the base case and the inductive step, the claim holds.

Since
\[
\Lambda^kX=\Lambda^kL\Lambda_1^{-k}\Lambda_1^kU=Q_k(\hat{R}_k\Lambda_1^kU)
\]
for some upper triangular \(\hat{R}_k\), the \(Q\)-factor in the QR factorization of \(\Lambda^kX\) is \(Q_k\). The claim gives
\[
|Q_k(i,j)|=\begin{cases}
|(L(1:m,:)^{-1})(j,i)|\dfrac{|\lambda_j|^k}{|\lambda_i|^k}(1+o(1)),\quad&i\in[j-1]\\
1+o(1),\quad&i=j\\
|L(i,j)|\dfrac{|\lambda_i|^k}{|\lambda_j|^k}(1+o(1)),\quad&i\in[n]\setminus[j]
\end{cases}.
\]
By the uniqueness of the QR factorization, different choices of the QR factorization do not affect \(|Q_k(i,j)|\), which completes the proof of Lemma~\ref{lemma:withoutRR}.

\section{Subspace iteration with the Rayleigh--Ritz procedure}\label{section:withRR}

\begin{lemma}\label{lemma:withRR}
Suppose \(\mathbb{F}=\mathbb{R}\) or \(\mathbb{C}\), \(n\), \(m\in\mathbb{N}_+\) with \(n>m\), and \(\Lambda=\mathrm{diag}(\lambda_1,\dots,\lambda_n)\in\mathbb{F}^{n\times n}\), with \(|\lambda_1|>\dots>|\lambda_n|>0\). Assume \(X\in\mathbb{F}^{n\times m}\) is such that \(X(1:m,:)\) is invertible and all entries of \(X(m+1:n,:)X(1:m,:)^{-1}\) are nonzero. 

Then there exists \(M_6=M_6(n,m,\Lambda,X)>0\), \(M_7=M_7(n,m,\Lambda,X)>0\), \(k_7=k_7(n,m,\Lambda,X)\in\mathbb{N}_+\) such that for every \(k\geq k_7\),

\begin{itemize}
\item The \(m\) Ritz values of \(\Lambda\) projected on \(\mathrm{span}(\Lambda^kX)\) have distinct modulus; denote them by \(\{\theta_j\}_{j\in[m]}\) where \(|\theta_1|>\dots>|\theta_m|\).
\item Let \(Z_k(:,j)\) be a unit Ritz vector corresponding to \(\theta_j\). Then for every \(j\in[m]\), we have
\begin{align*}
M_7\dfrac{|\lambda_{m+1}|^{2k}}{|\lambda_i|^k|\lambda_j|^k}&\leq|Z_k(i,j)|\leq M_6\dfrac{|\lambda_{m+1}|^{2k}}{|\lambda_i|^k|\lambda_j|^k},\quad i\in[m]\setminus\{j\};\\
M_7\dfrac{|\lambda_i|^k}{|\lambda_j|^k}&\leq|Z_k(i,j)|\leq M_6\dfrac{|\lambda_i|^k}{|\lambda_j|^k},\quad i\in[n]\setminus[m];\\
M_7\dfrac{|\lambda_{m+1}|^{2k}}{|\lambda_j|^{2k}}&\leq1-|Z_k(j,j)|\leq M_6\dfrac{|\lambda_{m+1}|^{2k}}{|\lambda_j|^{2k}}.
\end{align*}
\end{itemize}
\end{lemma}

\begin{theorem}
Suppose \(\mathbb{F}=\mathbb{R}\) or \(\mathbb{C}\), \(n\), \(m\in\mathbb{N}_+\) with \(n>m\), and \(\Lambda=\mathrm{diag}(\lambda_1,\dots,\lambda_n)\in\mathbb{F}^{n\times n}\), with \(|\lambda_1|>\dots>|\lambda_n|>0\). Assume \(V\in\mathbb{F}^{n\times n}\) satisfies \(VV^*=V^*V=I_n\). Define \(A:=V\Lambda V^*\)
. Assume \(Y\in\mathbb{F}^{n\times m}\). Define \(X:=V^*Y\), and assume \(X(1:m,:)\) is invertible and all entries of \(X(m+1:n,:)X(1:m,:)^{-1}\) are nonzero.

Then there exists \(M_6=M_6(n,m,\Lambda,X)>0\), \(M_7=M_7(n,m,\Lambda,X)>0\), \(k_7=k_7(n,m,\Lambda,X)\in\mathbb{N}_+\) such that for every \(k\geq k_7\),

\begin{itemize}
\item The \(m\) Ritz values of \(A\) projected on \(\mathrm{span}(A^kY)\) have distinct modulus; denote them by \(\{\theta_j\}_{j\in[m]}\) where \(|\theta_1|>\dots>|\theta_m|\).
\item Let \(W_k(:,j)\) be a unit Ritz vector corresponding to \(\theta_j\). Then for every \(j\in[m]\), we have
\begin{align*}
M_7\dfrac{|\lambda_{m+1}|^{2k}}{|\lambda_i|^k|\lambda_j|^k}&\leq|V(:,i)^*W_k(:,j)|\leq M_6\dfrac{|\lambda_{m+1}|^{2k}}{|\lambda_i|^k|\lambda_j|^k},\quad i\in[m]\setminus\{j\};\\
M_7\dfrac{|\lambda_i|^k}{|\lambda_j|^k}&\leq|V(:,i)^*W_k(:,j)|\leq M_6\dfrac{|\lambda_i|^k}{|\lambda_j|^k},\quad i\in[n]\setminus[m];\\
M_7\dfrac{|\lambda_{m+1}|^{2k}}{|\lambda_j|^{2k}}&\leq1-|V(:,j)^*W_k(:,j)|\leq M_6\dfrac{|\lambda_{m+1}|^{2k}}{|\lambda_j|^{2k}}.
\end{align*}
\end{itemize}
\end{theorem}
\begin{proof}
This follows by applying a change of orthonormal basis to the conclusion of the preceding lemma.
\end{proof}

We now prove Lemma~\ref{lemma:withRR}.

\subsection{Notation}

\begin{itemize}
\item \(\Lambda_1:=\mathrm{diag}(\lambda_1,\dots,\lambda_m)\), \(\Lambda_2:=\mathrm{diag}(\lambda_{m+1},\dots,\lambda_n)\).
\item \(\delta:=\min_{i\neq j}||\lambda_i|-|\lambda_j||>0\).
\item \(S=X(m+1:n,:)X(1:m,:)^{-1}\).
\item \(C:=\max\{1,|\lambda_{m+1}|\}\cdot\max_{p,q\in[m]}\{\sum_{r\in[n-m]}|S(r,p)S(r,q)|\}>0\).
\item \(M_0:=Cm(|\lambda_1|+2)\max\{1,2/\delta\}>0\).
\item \(k\in\mathbb{N}_+\).
\item \(H:=\Lambda_2^kS\Lambda_1^{-k}\), \(\Phi=\begin{bmatrix}
I_m\\
H
\end{bmatrix}\), \(E:=H^* H\), \(F:=H^*\Lambda_2H\).
\end{itemize}

\subsection{Auxiliary lemmas and propositions}

\begin{definition}
Set
\[
\rho_{p,q}=\dfrac{|\lambda_{m+1}|^2}{|\lambda_p||\lambda_q|},\quad p,q\in[m];\qquad\rho_{\max}=\dfrac{|\lambda_{m+1}|^2}{|\lambda_m|^2}<1.
\]
\end{definition}

\begin{proposition}
For \(p,q,r\in[m]\), we have
\[
\rho_{p,q}\rho_{q,r}\leq\rho_{\max}\rho_{p,r}.
\]
\end{proposition}
\begin{proof}
\[
\rho_{p,q}\rho_{q,r}=\dfrac{|\lambda_{m+1}|^2}{|\lambda_p||\lambda_q|}\dfrac{|\lambda_{m+1}|^2}{|\lambda_q||\lambda_r|}=\dfrac{|\lambda_{m+1}|^2}{|\lambda_q|^2}\dfrac{|\lambda_{m+1}|^2}{|\lambda_p||\lambda_r|}\leq\rho_{\max}\rho_{p,r}.
\]
\end{proof}

\begin{proposition}
\begin{enumerate}
\item \(H(p,q)=S(p,q)\dfrac{\lambda_{m+p}^k}{\lambda_{q}^k},\quad p\in[n-m],q\in[m]\).
\item \(E(p,q)=\sum_{r\in[n-m]}\overline{S(r,p)}S(r,q)\dfrac{\overline{\lambda_{m+r}}^k\lambda_{m+r}^k}{\overline{\lambda_p}^{k}\lambda_q^k},\quad p,q\in[m]\).
\item \(F(p,q)=\sum_{r\in[n-m]}\lambda_{m+r}\overline{S(r,p)}S(r,q)\dfrac{\overline{\lambda_{m+r}}^k\lambda_{m+r}^k}{\overline{\lambda_p}^{k}\lambda_q^k},\quad p,q\in[m]\).
\end{enumerate}
\end{proposition}

\begin{proposition}
\(|E(p,q)|\leq C\rho_{p,q}^k\), \(|F(p,q)|\leq C\rho_{p,q}^k\).
\end{proposition}
\begin{proof}
\[
|E(p,q)|\leq\sum_{r\in[n-m]}|S(r,p)S(r,q)|\dfrac{|\lambda_{m+r}|^{2k}}{|\lambda_p|^k|\lambda_q|^k}\leq\left(\sum_{r\in[n-m]}|S(r,p)S(r,q)|\right)\rho_{p,q}^k\leq C\rho_{p,q}^k.
\]
\[
|F(p,q)|\leq\sum_{r\in[n-m]}|\lambda_{m+r}S(r,p)S(r,q)|\dfrac{|\lambda_{m+r}|^{2k}}{|\lambda_p|^k|\lambda_q|^k}\leq\left(|\lambda_{m+1}|\sum_{r\in[n-m]}|S(r,p)S(r,q)|\right)\rho_{p,q}^k\leq C\rho_{p,q}^k.
\]
\end{proof}

\subsection{Construction of the metric space \((B,d)\) and the map \(T\)}

Fix \(j\in[m]\). Fix \(y(j)=1\), where \(y\in\mathbb{F}^m\), and let \(\hat{y}\in\mathbb{F}^{m-1}\) denote \(y\) with its \(j\)-th component removed. Consider the set
\[
B=\left\{(\theta,\hat{y})\in\mathbb{F}\times\mathbb{F}^{m-1}\middle|\dfrac{|\theta-\lambda_j|}{\rho_{j,j}^k}\leq M_0,\max_{q\neq j}\dfrac{|y(q)|}{\rho_{j,q}^k}\leq M_0\right\}.
\]

Define metric \(d\) on \(B\):
\[
d((\theta_a,\hat{y}_a),(\theta_b,\hat{y}_b)):=\max\left\{\dfrac{|\theta_a-\theta_b|}{\rho_{j,j}^k},\max_{q\neq j}\dfrac{|y_a(q)-y_b(q)|}{\rho_{j,q}^k}\right\}.
\]

\begin{proposition}
\(d\) is indeed a metric, and the metric space \((B,d)\) is complete.
\end{proposition}
\begin{proof}
The proof is straightforward.
\end{proof}

There exists \(k_1\) such that, for \(k\geq k_1\), \(M_0\rho_{\max}^k<\delta/2\).

For \(k\geq k_1\), define \(T:B\to\mathbb{F}\times\mathbb{F}^{m-1}\) by \((\theta,\hat{y})\mapsto(\theta',\hat{y}')\).
\begin{align*}
\theta'&=\lambda_j-\theta\sum_{q}E(j,q)y(q)+\sum_{q}F(j,q)y(q),\\
y'(p)&=\dfrac{1}{\lambda_p-\theta}\left[\theta\sum_{q}E(p,q)y(q)-\sum_{q}F(p,q)y(q)\right],\quad p\in[m]\setminus\{j\}.
\end{align*}

\subsection{\(T(B)\subseteq B\)}

There exists \(k_2\geq k_1\) such that, for \(k\geq k_2\), \(M_0\rho_{\max}^k<1\).

\begin{lemma}
If \(k\geq k_2\), then for every \(p\in[m]\), we have
\[
\left|\theta\sum_{q}E(p,q)y(q)-\sum_{q}F(p,q)y(q)\right|\leq(|\lambda_1|+2)mC\rho_{p,j}^k.
\]
\end{lemma}
\begin{proof}
\begin{align*}
&\left|\theta\sum_{q}E(p,q)y(q)-\sum_{q}F(p,q)y(q)\right|\\
&\leq|\theta|\sum_{q}|E(p,q)||y(q)|+\sum_{q}|F(p,q)||y(q)|\\
&\leq(|\lambda_j|+M_0\rho_{j,j}^k)\left(|E(p,j)|+\sum_{q\neq j}|E(p,q)||y(q)|\right)+|F(p,j)|+\sum_{q\neq j}|F(p,q)||y(q)|\\
&\leq(|\lambda_1|+2)\left(C\rho_{p,j}^k+\sum_{q\neq j}C\rho_{p,j}^kM_0\rho_{\max}^k\right)\\
&\leq(|\lambda_1|+2)mC\rho_{p,j}^k.
\end{align*}
\end{proof}

\begin{theorem}
If \(k\geq k_2\), then \(T(B)\subseteq B\).
\end{theorem}
\begin{proof}
\(\forall (\theta,\hat{y})\in B\). By the preceding lemma and the definition of \(M_0\),
\[
|\theta'-\lambda_j|=\left|-\theta\sum_{q}E(j,q)y(q)+\sum_{q}F(j,q)y(q)\right|\leq (|\lambda_1|+2)mC\rho_{j,j}^k\leq M_0\rho_{j,j}^k.
\]
For every \(p\in[m]\setminus\{j\}\),
\[
|\lambda_p-\theta|\geq|\lambda_p-\lambda_j|-|\lambda_j-\theta|\geq\delta-\delta/2=\delta/2,
\]
by the preceding lemma and the definition of \(M_0\),
\[
|y'(p)|=\dfrac{1}{|\lambda_p-\theta|}\left|\theta\sum_{q}E(p,q)y(q)-\sum_{q}F(p,q)y(q)\right|\leq\dfrac{2}{\delta}(|\lambda_1|+2)mC\rho_{p,j}^k\leq M_0\rho_{p,j}^k.
\]
Hence we have \(T(\theta,\hat{y})=(\theta',\hat{y}')\in B\). Therefore \(T(B)\subseteq B\).
\end{proof}

\subsection{\(T\) is a strict contraction}

\begin{lemma}
If \(k\geq k_2\), then for all \((\theta_a,\hat{y}_a)\), \((\theta_b,\hat{y}_b)\in B\) and every \(p\in[m]\), we have
\begin{align*}
&\left|\left[\theta_a\sum_{q}E(p,q)y_a(q)-\sum_{q}F(p,q)y_a(q)\right]-\left[\theta_b\sum_{q}E(p,q)y_b(q)-\sum_{q}F(p,q)y_b(q)\right]\right|\\
&\leq Cd(a,b)\rho_{p,j}^k\rho_{\max}^km(|\lambda_1|+3).
\end{align*}
\end{lemma}
\begin{proof}
\begin{align*}
&\left|\left[\theta_a\sum_{q}E(p,q)y_a(q)-\sum_{q}F(p,q)y_a(q)\right]-\left[\theta_b\sum_{q}E(p,q)y_b(q)-\sum_{q}F(p,q)y_b(q)\right]\right|\\
&\leq\left|\sum_{q}E(p,q)\theta_ay_a(q)-\sum_{q}E(p,q)\theta_by_b(q)\right|+\left|\sum_{q}F(p,q)(y_a(q)-y_b(q))\right|\\
&\leq\sum_{q}|E(p,q)||\theta_ay_a(q)-\theta_by_b(q)|+\sum_{q}|F(p,q)||y_a(q)-y_b(q)|\\
&=|E(p,j)||\theta_a-\theta_b|+\sum_{q\neq j}|E(p,q)||\theta_ay_a(q)-\theta_by_b(q)|+\sum_{q\neq j}|F(p,q)||y_a(q)-y_b(q)|\\
&\leq C\rho_{p,j}^k\rho_{j,j}^kd(a,b)+\sum_{q\neq j}C\rho_{p,q}^k(|\theta_a||y_a(q)-y_b(q)|+|\theta_a-\theta_b||y_b(q)|)+\sum_{q\neq j}C\rho_{p,q}^k\rho_{j,q}^kd(a,b)\\
&\leq mC\rho_{p,j}^k\rho_{\max}^kd(a,b)+\sum_{q\neq j}C\rho_{p,q}^k((|\lambda_1|+1)\rho_{j,q}^kd(a,b)+\rho_{j,j}^kd(a,b))\\
&\leq Cd(a,b)\rho_{p,j}^k\rho_{\max}^km(|\lambda_1|+3).
\end{align*}
\end{proof}

There exists \(k_3\geq k_2\) such that, for \(k\geq k_3\), \(C\rho_{\max}^km(|\lambda_1|+3)\leq1/2\).

\begin{proposition}
If \(k\geq k_3\), then for all \((\theta_a,\hat{y}_a)\), \((\theta_b,\hat{y}_b)\in B\), we have \(\dfrac{|\theta_a'-\theta_b'|}{\rho_{j,j}^k}\leq\dfrac{1}{2}d(a,b)\).
\end{proposition}
\begin{proof}
\(\forall(\theta_a,\hat{y}_a)\), \((\theta_b,\hat{y}_b)\in B\). Set \((\theta_a',\hat{y}_a')=T(\theta_a,\hat{y}_a)\), \((\theta_b',\hat{y}_b')=T(\theta_b,\hat{y}_b)\). Then

\begin{align*}
&|\theta_a'-\theta_b'|\\
&=\left|\left(-\theta_a\sum_{q}E(j,q)y_a(q)+\sum_{q}F(j,q)y_a(q)\right)-\left(-\theta_b\sum_{q}E(j,q)y_b(q)+\sum_{q}F(j,q)y_b(q)\right)\right|\\
&\leq Cd(a,b)\rho_{j,j}^k\rho_{\max}^km(|\lambda_1|+3)\leq\dfrac{1}{2}\rho_{j,j}^kd(a,b),
\end{align*}

Hence
\[
\dfrac{|\theta_a'-\theta_b'|}{\rho_{j,j}^k}\leq\dfrac{1}{2}d(a,b).
\]
\end{proof}

There exists \(k_4\geq k_3\) such that, for \(k\geq k_4\), \(C\dfrac{8}{\delta^2}\rho_{\max}^km(|\lambda_1|+3)^2\leq1/2\).

\begin{proposition}
If \(k\geq k_4\), then for all \((\theta_a,\hat{y}_a)\), \((\theta_b,\hat{y}_b)\in B\) and every \(p\in[m]\setminus\{j\}\), we have \(\dfrac{|y_a'(p)-y_b'(p)|}{\rho_{p,j}^k}\leq\dfrac{1}{2}d(a,b)\).
\end{proposition}
\begin{align*}
&|y_a'(p)-y_b'(p)|\\
&=\left|\dfrac{1}{\lambda_p-\theta_a}\left[\theta_a\sum_{q}E(p,q)y_a(q)-\sum_{q}F(p,q)y_a(q)\right]-\dfrac{1}{\lambda_p-\theta_b}\left[\theta_b\sum_{q}E(p,q)y_b(q)-\sum_{q}F(p,q)y_b(q)\right]\right|\\
&\leq\dfrac{4}{\delta^2}\left|(\lambda_p-\theta_b)\left[\theta_a\sum_{q}E(p,q)y_a(q)-\sum_{q}F(p,q)y_a(q)\right]-(\lambda_p-\theta_a)\left[\theta_b\sum_{q}E(p,q)y_b(q)-\sum_{q}F(p,q)y_b(q)\right]\right|\\
&\leq\dfrac{4}{\delta^2}|\theta_a-\theta_b|\left|\theta_a\sum_{q}E(p,q)y_a(q)-\sum_{q}F(p,q)y_a(q)\right|\\
&\quad+\dfrac{4}{\delta^2}|\lambda_p-\theta_a|\left|\left[\theta_a\sum_{q}E(p,q)y_a(q)-\sum_{q}F(p,q)y_a(q)\right]-\left[\theta_b\sum_{q}E(p,q)y_b(q)-\sum_{q}F(p,q)y_b(q)\right]\right|\\
&\leq\dfrac{4}{\delta^2}\rho_{j,j}^kd(a,b)(|\lambda_1|+2)mC\rho_{p,j}^k+\dfrac{4}{\delta^2}(2|\lambda_1|+1)Cd(a,b)\rho_{p,j}^k\rho_{\max}^km(|\lambda_1|+3)\\
&\leq\rho_{p,j}^kd(a,b)C\dfrac{8}{\delta^2}\rho_{\max}^km(|\lambda_1|+3)^2\\
&\leq\dfrac{1}{2}\rho_{p,j}^kd(a,b).
\end{align*}

So
\[
\dfrac{|y_a'(p)-y_b'(p)|}{\rho_{p,j}^k}\leq\dfrac{1}{2}d(a,b).
\]

\begin{theorem}
If \(k\geq k_4\), then for all \((\theta_a,\hat{y}_a)\), \((\theta_b,\hat{y}_b)\in B\), we have
\(d(a',b')\leq\dfrac{1}{2}d(a,b)\).
\end{theorem}
\begin{proof}
\[
d(a',b')=\max\left\{\dfrac{|\theta_a'-\theta_b'|}{\rho_{j,j}^k},\max_{q\neq j}\dfrac{|y_a'(q)-y_b'(q)|}{\rho_{j,q}^k}\right\}\leq\dfrac{1}{2}d(a,b).
\]
\end{proof}

\subsection{Application of the Banach fixed point theorem}

So far we have proved that if \(k\geq k_4\), then the following hold:
\begin{enumerate}
\item \((B,d)\) is a nonempty complete metric space;
\item \(T(B)\subseteq B\);
\item \(\forall a,b\in B\), \(d(T(a),T(b))\leq\dfrac{1}{2}d(a,b)\).
\end{enumerate}
By the Banach fixed point theorem, \(T\) has a unique fixed point in \(B\). Denote it by \((\theta_j,\hat{y}_j)\). By \(T(\theta_j,\hat{y}_j)=(\theta_j,\hat{y}_j)\), we know that\((\Lambda_1-\theta_jI_m)y_j=(\theta_jE-F)y_j\), which gives \((\Lambda_1+F)y_j=(I_m+E)y_j\theta_j\), and hence \(\Phi^*\Lambda\Phi y_j=\Phi^*\Phi y_j\theta_j\). Note that
\[
\Phi=\begin{bmatrix}
I_m\\
H
\end{bmatrix}=\begin{bmatrix}
I_m\\
\Lambda_2^kS\Lambda_1^{-k}
\end{bmatrix}=\begin{bmatrix}
\Lambda_1^kX(1:m,:)\\
\Lambda_2^kX(m+1:n,:)
\end{bmatrix}X(1:m,:)^{-1}\Lambda_1^{-k}=\Lambda^kXX(1:m,:)^{-1}\Lambda_1^{-k},
\]
thus, \((\theta_j,\Phi y_j)\) is a Ritz pair of \(\Lambda\) projected on \(\mathrm{span}(\Lambda^kX)\).

\((\theta_j,\hat{y}_j)\in B\) implies \(|\theta_j-\lambda_j|\leq M_0\rho_{\max}^k<\delta/2\). Therefore, for every \(j\in[m-1]\),
\[
|\theta_j|>|\lambda_j|-\delta/2=|\lambda_{j+1}|+|\lambda_j|-|\lambda_{j+1}|-\delta/2\geq |\lambda_{j+1}|+\delta/2>|\theta_{j+1}|.
\]
Therefore \(\{(\theta_j,\Phi y_j)\}_{j\in[m]}\) are exactly the \(m\) Ritz pairs of \(\Lambda\) projected on \(\mathrm{span}(\Lambda^kX)\), with \(|\theta_1|>\dots>|\theta_m|\).

Since \((\theta_j,\hat{y}_j)\in B\), we have the following proposition.

\begin{proposition}
(Upper bound for \(|y_j(q)|\)) If \(k\geq k_4\), then for every \(j\in[m]\) and every \(q\in[m]\setminus\{j\}\), we have
\[
|y_j(q)|\leq M_0\rho_{j,q}^k.
\]
\end{proposition}

\subsection{Lower bound for \(|y_j(q)|\)}

There exists \(k_5\geq k_4\) such that, for \(k\geq k_5\), \(\sum_{r=2}^{n-m}(2|\lambda_1|+1)|S(r,p)S(r,j)|\dfrac{|\lambda_{m+2}|^{2k}}{|\lambda_{m+1}|^{2k}}\leq\dfrac{\delta}{4}|S(1,p)S(1,j)|\) and \((|\lambda_1|+2)mCM_0\rho_{\max}^k\leq\dfrac{\delta}{8}|S(1,p)S(1,j)|\) for all \(p,j\in[m]\).

Set \(M_1=\min_{p,j\in[m]}\left\{\dfrac{\delta|S(1,p)S(1,j)|}{(2|\lambda_1|+1)8}\right\}\).

\begin{proposition}
If \(k\geq k_5\), then for every \(j\in[m]\) and every \(p\in[m]\setminus\{j\}\), we have \(|y_j(p)|\geq M_1\rho_{p,j}^k\).
\end{proposition}
\begin{proof}
We know that
\begin{align*}
&|\theta_jE(p,j)-F(p,j)|\\
&=\left|\sum_{r\in[n-m]}\theta_j\overline{S(r,p)}S(r,j)\dfrac{\overline{\lambda_{m+r}}^k\lambda_{m+r}^k}{\overline{\lambda_p}^k\lambda_j^k}-\sum_{r\in[n-m]}\lambda_{m+r}\overline{S(r,p)}S(r,j)\dfrac{\overline{\lambda_{m+r}}^k\lambda_{m+r}^k}{\overline{\lambda_p}^k\lambda_j^k}\right|\\
&\geq|\theta_j-\lambda_{m+1}||S(1,p)S(1,j)|\rho_{p,j}^k-\sum_{r=2}^{n-m}(2|\lambda_1|+1)|S(r,p)S(r,j)|\rho_{p,j}^k\dfrac{|\lambda_{m+2}|^{2k}}{|\lambda_{m+1}|^{2k}}\\
&\geq\dfrac{\delta}{2}|S(1,p)S(1,j)|\rho_{p,j}^k-\dfrac{\delta}{4}|S(1,p)S(1,j)|\rho_{p,j}^k\\
&=\dfrac{\delta}{4}|S(1,p)S(1,j)|\rho_{p,j}^k.
\end{align*}

Therefore

\begin{align*}
|y_j(p)|&=\dfrac{1}{|\lambda_p-\theta_j|}\left|\theta_j\sum_{q}E(p,q)y_j(q)-\sum_{q}F(p,q)y_j(q)\right|\\
&\geq\dfrac{1}{2|\lambda_1|+1}\left(|\theta_jE(p,j)-F(p,j)|-|\theta_j|\sum_{q\neq j}|E(p,q)||y_j(q)|-\sum_{q\neq j}|F(p,q)||y_j(q)|\right)\\
&\geq\dfrac{1}{2|\lambda_1|+1}\left(|\theta_jE(p,j)-F(p,j)|-(|\lambda_1|+2)\sum_{q\neq j}C\rho_{p,q}^kM_0\rho_{j,q}^k\right)\\
&\geq\dfrac{1}{2|\lambda_1|+1}\left(\dfrac{\delta}{4}|S(1,p)S(1,j)|\rho_{p,j}^k-(|\lambda_1|+2)mC\rho_{p,j}^kM_0\rho_{\max}^k\right)\\
&\geq\dfrac{\delta|S(1,p)S(1,j)|}{(2|\lambda_1|+1)8}\rho_{p,j}^k\\
&\geq M_1\rho_{p,j}^k.
\end{align*}
\end{proof}

\subsection{Returning to Ritz vectors}

\(\forall j\in[m]\). The \(j\)-th Ritz vector is
\[
\Phi y_j=\begin{bmatrix}
I_m\\
H
\end{bmatrix}y_j=\begin{bmatrix}
y_j\\
Hy_j
\end{bmatrix}
\]

For every \(q\in[m]\), we have
\[
(\Phi y_j)(q)=y_j(q).
\]

For every \(p\in[n-m]\), we have
\[
(\Phi y_j)(m+p)=(Hy_j)(p)=\sum_{q}H(p,q)y_j(q)=H(p,j)+\sum_{q\neq j}H(p,q)y_j(q).
\]

There exists \(k_6\geq k_5\) such that, for \(k\geq k_6\), \(\sum_{q\neq j}|S(p,q)|M_0\rho_{\max}^k\leq\dfrac{1}{2}|S(p,j)|\) for all \(p\in[n-m]\), \(j\in[m]\).

Set \(M_2=\dfrac{3}{2}\max_{p\in[n-m],j\in[m]}\{|S(p,j)|\}\) and \(M_3=\dfrac{1}{2}\min_{p\in[n-m],j\in[m]}\{|S(p,j)|\}\).

Set \(M_4=\max\{M_0,M_2\}\) and \(M_5=\min\{M_1,M_3\}\).

\begin{proposition}
If \(k\geq k_6\), then for every \(j\in[m]\), the following hold:
\begin{itemize}
\item \((\Phi y_j)(j)=1\);
\item \(M_5\dfrac{|\lambda_{m+1}|^{2k}}{|\lambda_j|^k|\lambda_q|^k}\leq|(\Phi y_j)(q)|\leq M_4\dfrac{|\lambda_{m+1}|^{2k}}{|\lambda_j|^k|\lambda_q|^k},\quad\forall q\in[m]\setminus\{j\}\);
\item \(M_5\dfrac{|\lambda_q|^k}{|\lambda_j|^k}\leq|(\Phi y_j)(q)|\leq M_4\dfrac{|\lambda_q|^k}{|\lambda_j|^k},\quad\forall q\in[n]\setminus[m]\).
\end{itemize}
\end{proposition}
\begin{proof}
It suffices to prove the case \(q\in[n]\setminus[m]\).

\(\forall p\in[n-m]\),
\begin{align*}
&|(\Phi y_j)(m+p)|\leq|H(p,j)|+\sum_{q\neq j}|H(p,q)||y_j(q)|\\
&\leq |S(p,j)|\dfrac{|\lambda_{m+p}|^k}{|\lambda_j|^k}+\sum_{q\neq j}|S(p,q)|\dfrac{|\lambda_{m+p}|^k}{|\lambda_j|^k}M_0\rho_{\max}^k\\
&\leq\dfrac{3}{2}|S(p,j)|\dfrac{|\lambda_{m+p}|^k}{|\lambda_j|^k}\leq M_2\dfrac{|\lambda_{m+p}|^k}{|\lambda_j|^k}\leq M_4\dfrac{|\lambda_{m+p}|^k}{|\lambda_j|^k}.
\end{align*}
Similarly we have
\[
|(\Phi y_j)(m+p)|\geq|H(p,j)|-\sum_{q\neq j}|H(p,q)||y_j(q)|\geq M_5\dfrac{|\lambda_{m+p}|^k}{|\lambda_j|^k}.
\]
\end{proof}

\subsection{Normalization in the 2-norm}

The true Ritz pairs are given by \(\left\{\left(\theta_j,\dfrac{\Phi y_j}{\|\Phi y_j\|_2}\right)\right\}_{j\in[m]}\). Define \(z_j=\dfrac{\Phi y_j}{\|\Phi y_j\|_2}\).

There exists \(k_7\geq k_6\) such that, for \(k\geq k_7\), \(nM_4^2\dfrac{|\lambda_{m+1}|^{2k}}{|\lambda_m|^{2k}}\leq 3\). Set \(M_6=\max\left\{M_4,nM_4^2\right\}\) and \(M_7=\min\left\{\dfrac{1}{2}M_5,\dfrac{1}{8}M_5^2\right\}\).

\begin{proposition}
Suppose \(k\geq k_7\). For every \(j\in[m]\), we have
\begin{align*}
M_7\dfrac{|\lambda_{m+1}|^{2k}}{|\lambda_j|^k|\lambda_q|^k}&\leq|z_j(q)|\leq M_6\dfrac{|\lambda_{m+1}|^{2k}}{|\lambda_j|^k|\lambda_q|^k},\quad\forall q\in[m]\setminus\{j\};\\
M_7\dfrac{|\lambda_q|^k}{|\lambda_j|^k}&\leq|z_j(q)|\leq M_6\dfrac{|\lambda_q|^k}{|\lambda_j|^k},\quad\forall q\in[n]\setminus[m];\\
M_7\dfrac{|\lambda_{m+1}|^{2k}}{|\lambda_j|^{2k}}&\leq1-|z_j(j)|\leq M_6\dfrac{|\lambda_{m+1}|^{2k}}{|\lambda_j|^{2k}}.
\end{align*}
\end{proposition}
\begin{proof}
\(\forall j\in[m]\).
\begin{enumerate}
\item For every \(q\in[m]\setminus\{j\}\), \(|(\Phi y_j)(q)|\leq M_4\dfrac{|\lambda_{m+1}|^{2k}}{|\lambda_j|^k|\lambda_q|^k}\leq M_4\dfrac{|\lambda_{m+1}|^k}{|\lambda_m|^k}\),
\item For every \(q\in[n]\setminus[m]\), \(|(\Phi y_j)(q)|\leq M_4\dfrac{|\lambda_q|^k}{|\lambda_j|^k}\leq M_4\dfrac{|\lambda_{m+1}|^k}{|\lambda_m|^k}\).
\end{enumerate}

Since \(k\geq k_7\),
\[
1=|(\Phi y_j)(j)|^2\leq\|\Phi y_j\|_2^2\leq 1+nM_4^2\dfrac{|\lambda_{m+1}|^{2k}}{|\lambda_m|^{2k}}\leq 4.
\]
Therefore \(1\leq\|\Phi y_j\|_2\leq 2\).

Thus, for \(q\in[m]\setminus\{j\}\),
\[
M_7\dfrac{|\lambda_{m+1}|^{2k}}{|\lambda_j|^k|\lambda_q|^k}\leq\dfrac{1}{2}M_5\dfrac{|\lambda_{m+1}|^{2k}}{|\lambda_j|^k|\lambda_q|^k}\leq|z_j(q)|\equiv\dfrac{|(\Phi y_j)(q)|}{\|\Phi y_j\|_2}\leq M_4\dfrac{|\lambda_{m+1}|^{2k}}{|\lambda_j|^k|\lambda_q|^k}\leq M_6\dfrac{|\lambda_{m+1}|^{2k}}{|\lambda_j|^k|\lambda_q|^k}.
\]
for \(q\in[n]\setminus[m]\),
\[
M_7\dfrac{|\lambda_q|^k}{|\lambda_j|^k}\leq\dfrac{1}{2}M_5\dfrac{|\lambda_q|^k}{|\lambda_j|^k}\leq|z_j(q)|\equiv\dfrac{|(\Phi y_j)(q)|}{\|\Phi y_j\|_2}\leq M_4\dfrac{|\lambda_q|^k}{|\lambda_j|^k}\leq M_6\dfrac{|\lambda_q|^k}{|\lambda_j|^k}.
\]
for \(q=j\), using \(1-x\leq\sqrt{1-x}\leq1-x/2\) for \(x\in[0,1]\), we have
\[
1-\sum_{q\in[n],q\neq j}|z_j(q)|^2\leq|z_j(j)|\equiv\sqrt{1-\sum_{q\in[n],q\neq j}|z_j(q)|^2}\leq 1-\dfrac{1}{2}\sum_{q\in[n],q\neq j}|z_j(q)|^2,
\]
and hence
\begin{align*}
&M_7\dfrac{|\lambda_{m+1}|^{2k}}{|\lambda_j|^{2k}}\leq\dfrac{1}{8}M_5^2\dfrac{|\lambda_{m+1}|^{2k}}{|\lambda_j|^{2k}}\leq\dfrac{1}{2}\sum_{q\in[n],q\neq j}|z_j(q)|^2\leq1-|z_j(j)|\\
&\leq\sum_{q\in[n],q\neq j}|z_j(q)|^2\leq nM_4^2\dfrac{|\lambda_{m+1}|^{2k}}{|\lambda_j|^{2k}}\leq M_6\dfrac{|\lambda_{m+1}|^{2k}}{|\lambda_j|^{2k}}.
\end{align*}
\end{proof}

All unit Ritz vectors corresponding to \(\theta_j\) are of the form \(z_je^{\imath\alpha}\) for \(\alpha\in\mathbb{R}\). It is clear that \(|(z_je^{\imath\alpha})(q)|\) is independent of \(\alpha\). This completes the proof of Lemma~\ref{lemma:withRR}.

\section{Conclusion}\label{section:conclusion}

In this paper, we have studied the element-wise convergence behavior of subspace iteration for a unitarily diagonalizable matrix
\[
A = V \Lambda V^*, \qquad \Lambda = \operatorname{diag}(\lambda_1,\dots,\lambda_n),
\]
with \(|\lambda_1| > \cdots > |\lambda_n| > 0\). All results are established for both \(\mathbb{F} = \mathbb{R}\) and \(\mathbb{F} = \mathbb{C}\). All element-wise statements below are understood in the eigenvector coordinate system associated with \(V\), that is, we consider the entries of \(V^*\tilde Q_k\) and \(V^*W_k\), rather than the entries of \(\tilde Q_k\) and \(W_k\) in the standard basis.

For subspace iteration without the Rayleigh--Ritz procedure, let \(\tilde Q_k \tilde R_k = A^k Y\) be the QR factorization. Under mild assumptions of \(Y\), the entries of \(V^*\tilde Q_k\) satisfy
\[
|V(:,i)^*\tilde Q_k(:,j)|=\Theta \left(\frac{|\lambda_{\max\{i,j\}}|^k}{|\lambda_{\min\{i,j\}}|^k}\right), \quad (i,j)\in[n]\times [m],\quad i\neq j,
\]
\[
1-|V(:,j)^*\tilde Q_k(:,j)|=\Theta\left( \max\left\{\frac{|\lambda_j|}{|\lambda_{j-1}|}, \frac{|\lambda_{j+1}|}{|\lambda_j|}\right\}^{2k}\right),\quad j\in[m].
\]

For subspace iteration with the Rayleigh--Ritz procedure, let \(W_k\) denote the Ritz vectors of \(A\) projected on \(\mathrm{span}(A^kY)\), with Ritz values sorted in the descending order by magnitude. Under mild assumptions of \(Y\), the entries of \(V^*W_k\) satisfy
\[
|V(:,i)^*W_k(:,j)|=\Theta \left(\frac{|\lambda_i|^k}{|\lambda_j|^k}\right),\quad (i,j)\in([n]\setminus[m])\times[m],
\]
\[
|V(:,i)^*W_k(:,j)|=\Theta \left(\frac{|\lambda_{m+1}|^{2k}}{|\lambda_i|^k|\lambda_j|^k}\right), \quad (i,j)\in[m]\times[m],\quad i\neq j,
\]
and
\[
1-|V(:,j)^*W_k(:,j)|=\Theta \left(\frac{|\lambda_{m+1}|^{2k}}{|\lambda_j|^{2k}}\right), \quad j\in[m].
\]
These results provide an element-wise description of the convergence of subspace iteration, both without and with the Rayleigh--Ritz procedure, in the eigenvector coordinate system.

A natural direction for future work is to extend the present element-wise theory from unitarily diagonalizable matrices to unitarily upper triangularizable matrices, and this generalization is left for future research.

\bibliographystyle{plain}
\bibliography{references}

\end{document}
